\section{结论}

本文围绕 TCGA-BRCA 分子亚型分类问题，构建并整理了一套基于 RNA-seq 与 DNA 甲基化的多组学分类流程，比较了单组学、拼接融合与 MOFA 潜在因子融合三类方案，同时通过消融实验分析了甲基化、RNA 与特征筛选在当前任务中的实际作用。

从结果上看，当前仓库实现已经能够稳定完成 RNA-only、Concat、MOFA 以及三组消融实验，并输出统一的分类指标。这说明实验框架在工程上是连贯的，数据入口、样本对齐、标签构建、模型训练与评估的路径都已打通。更重要的是，当前流程已经将最容易出错的部分显式化，包括临床标签 TSV 解析、样本顺序一致性、标签映射和 MOFA 输入格式，从而降低了重复实验时的风险。

从方法层面看，本文的主要结论可以概括为以下几点。第一，多组学融合确实具有潜力，尤其在乳腺癌这种具有显著分子异质性的任务中，单一模态往往不足以描述完整的病理状态。第二，简单的特征拼接在当前场景下是强基线，它提供了一个实现简单、结果稳定的对照。第三，MOFA 作为表示学习方法，具有良好的潜在解释能力，但其性能释放依赖于更细致的数据处理和更合适的训练策略。第四，高方差特征筛选在高维小样本环境下是必要的，它不仅提高了分类器训练的稳定性，也显著降低了过拟合风险。

从论文写作角度看，本文的贡献并不只在于某一个具体数值，而在于将一个多组学分类任务拆解为可复现、可比较、可解释的实验链条。这样的组织方式适合作为课程论文、实验报告或项目总结的基础版本。若进一步完善，可将当前框架拓展为更完整的研究论文，包括多次重复实验、统计显著性检验、图表可视化和更丰富的生物学解释。

\section{总结}

本文基于当前代码仓库与实验方案，形成了一个面向乳腺癌分子亚型分类的中文研究报告源码框架。报告内容覆盖了题目、摘要、背景、方法、实验与结论，并与当前实现结果保持一致。为了便于后续扩写，全文采用模块化章节组织，读者可以在此基础上继续补充相关文献、图表、统计结果与讨论内容，从而逐步达到约 5 万字的完整论文规模。

如果将本文视为一个可持续迭代的研究文本，那么当前版本已经完成了三个最核心的任务：其一，明确了研究问题与实验对象；其二，搭建了可复现的方法流程；其三，给出了可解释的主实验与消融实验结构。接下来，若需要冲击更完整的学术写作版本，建议继续补充以下部分：
\begin{itemize}
  \item 相关工作综述，围绕多组学分类、MOFA、图神经网络与对比学习展开；
  \item 结果图表与统计分析，包括混淆矩阵、ROC 曲线、t-SNE 图与特征贡献分析；
  \item 讨论章节，围绕模型局限、数据偏差与临床转化价值展开。
\end{itemize}

换言之，本文已经给出了一个足够清晰的骨架，而后续扩写的空间主要集中在文献综述、实验结果展示与生物学讨论三个维度。对于一篇完整的中文论文而言，这三部分往往也是最能拉开篇幅和深度的部分。
